\section{Lean Verification}\label{app:lean}

We verify selected results in Lean~4~\cite{lean4}, using mathlib~\cite{mathlib}. The artifact includes the proof sources, a theorem map, a complete manuscript coverage inventory, and reproducible compilation and axiom audits. It distinguishes complete results from specializations and conditional steps.

The general collision lemma is verified for arbitrary features and encoders and every $K\ge4$, without differentiability of $\psi$. The operational score formula \eqref{eq:exactscore} and its shrinking-pair clause are also verified for every nondegenerate two-feature encoder and $K\ge2$, under the stated feature curvature and endpoint regularity assumptions. The formal score is the infimum over actual admissible input pairs. The proof uses explicit two-input perturbations and one-sided endpoint derivatives.

For variance, the score identity additionally covers degenerate quadratic curves. Lean verifies the operational optimum over every quadratic coefficient curve at all positive powers:
$$
 \max_{\substack{\enc(s)=c+us+vs^2\\|\enc(s)|^2\le P\;\forall s\in I}}
 \Lambda_{\sq}(\enc;V)
 =\frac{4K^2P}{9+6\sqrt3},\qquad K\ge2.
$$
The complete quadratic and quartic Gram-objective optimizations, including continuum peak certificates and arbitrary positive slope parameter and power, are also checked. These coefficient-class theorems do not yet extend their competitors to every continuous encoder.

For encoder design, the artifact verifies the Schur-complement reduction, rank-one score constraints, and explicit planar reconstruction. The reconstruction is connected to the general operational score: a positive feasible SDP point yields a continuous, exact encoder satisfying the full peak constraint and the promised separation bound. The proof derives spanning anchors and a positive feasible point from the feature assumptions, establishes compactness, and proves maximum-separation attainment and its operational realization within the affine-feature class.

The rigidity development constructs genuine local moment inverses and neighborhoods, derives the derivative-kernel relations from local fiber constancy, integrates through both source endpoints, and proves coefficient independence from exactness. Its final representation theorem still assumes interior encoder differentiability and local moment-fiber constancy. Deriving these properties from continuity and exactness requires the remaining topological factorization and smoothing arguments; planar invariance of domain is not formalized here.

The decoder development verifies the sharp scalar amplification principle, including the zero-score impossibility result, and the explicit complex-channel variance decoder with corrected strip projection. It also checks coherent implementation-error bounds and peak-safe contraction before quantization. The small-noise development proves the derivative limit, integrable domination, dominated convergence in noise-power coordinates, and independent isotropic second-moment evaluation. It derives local differentiability of the actual variance decoder and its leading-MSE limit. The complete general-decoder gradient-to-trace specialization in Proposition~\ref{prop:general-small-noise} remains unfinished.

For the probabilistic design problem, Lean verifies convexity of the continuum Gram constraints, strict convexity of inverse-Gram risk, and convexity and monotonicity of the actual infimum frontier. The artifact also proves existence of a risk-minimizing Gram design at every floor from zero through the positive maximum, supplying the optimizer witnesses used by its uniqueness, plateau, activity, and strict-increase results. Endpoint continuity, the five-variable inverse-form equivalence, and the complete probabilistic and continuous-encoder links remain unfinished.

The artifact pins Lean~4.24.0 and a matching mathlib revision. Its audited declarations use only \texttt{propext}, \texttt{Classical.choice}, and \texttt{Quot.sound}, with no unfinished proofs or additional axioms. The unrestricted continuous-encoder theorem, complete probabilistic frontier, certified-grid analysis, and remaining appendix results are not fully formalized. The artifact therefore supports a claim about selected verified results, not all proofs in the paper.
